% verification/cards/<card-ID>.tex として保存し，verification.tex から \input する．
% \usepackage{verification} が必要．card-IDはファイル名と一致させる．
% V・Mは該当するものだけ残す（複数可）．不要な詳細欄は削除．
\begin{vcard}{E-semantic-title}{target-label (file.tex @ abc1234), scope: ...}{検証対象の短い題名}
% 未コミットの変更を含むなら保存版や差分も特定する．
\vwhat{V1 / V2 / V3 / V4 / V5: 具体的に調べる問い}
\vhow{M1 / M2 / M3 / M4: 実施した方法}
% 予定と実施済みを区別する．
\vresult{未確認: 判定範囲・理由}
% 確認済み / 部分的に確認済み / 問題あり / 未確認．確認範囲を超えない．
\vpaper
\begin{paperquote}
[PASTE VERBATIM PAPER LINES HERE]
\end{paperquote}
\vevidence
% 実施した作業・根拠，または詳細欄・コード・報告への参照．重複転記しない．
% 外部の形式検証結果を参照する場合：保存先・報告・成果物の版・対象宣言．
% 原稿の命題との対応，検証範囲，追加仮定・axiom・sorry等，未形式化部分の
% 確認済み情報と未確認情報を区別．参照元の検査報告と本SkillでのM1による照合も区別．
% 形式化・カーネル検査は別Skillの責務．Notes自体は形式検証と同等の保証を与えない．
\vstep
% 任意．非自明な推論の展開や別証明と証明できた範囲．
% \vuse
% 仮定の使用箇所が必要な場合．
\vlocator
% 依存の適用を調べる場合：内部labelまたは文献ID・著者・題名／書名・年／版・定理／頁と確認範囲．
% 外部文献にはDOI／URL等，入手済みならpapers/またはbook-pages/内のパスも記す．
\vsource
% 必要な主張の逐語引用．すぐ追える内部参照は保存版への参照でも可．
% 原文未入手の場合のみ \vtodo と必要な確認を記す．
% \vdict
% 既知の結果を本書の記号で導出・再定式化する場合は必須．原典と本書の記号対応．
\vgap
% 依存先が与える結論と必要な結論との差，埋める議論・障害．
\vdeps
% 任意．前提のlabel・card ID・文献．単なる参考と区別．
% 外部依存の確認状況・追跡の境界と理由，別証明との共通依存．
% \vinterface
% 接続が複雑な場合だけ．要求条件／引用元の結論／対象での照合／未確認事項．
\vtest
% 任意．実験・探索の方法・範囲・条件・結果・限界と再現用のコード等への参照．
% 価値のある失敗は障害・使用した仮定・再検討条件も記す．
\vopen
% 未確認箇所と次の確認．対象範囲外と現在のcorrectnessに関わる問題を区別．
% 後者はQ-semantic-issue等の項目IDを明記し，verification/to_verify.mdと同期．
% 解決後は項目ID・版・根拠をLogに残して両方の未解決欄から除く．
% \vproblem
% 任意．元の主張の検証と区別する探索課題．
\vlog
\begin{itemize}
\item YYYY-MM-DD: 作業・発見・変更・判断の根拠．
% 独立再検証時：担当・方法，版・範囲，渡した情報・先の判定を知っていたか，
% 根拠・不一致・共有依存・限界（報告への参照可）．
% 改訂時は旧版の根拠を保持し，現行版のResultと主要定理一覧を更新．
\end{itemize}
\end{vcard}
